% !TeX program = xelatex
% !TeX encoding = UTF-8
\documentclass{JXUSTmodeling}
\geometry{reset,a4paper,left=2.5cm,right=2.5cm,top=2.5cm,bottom=2.5cm}
\usepackage{float}
\graphicspath{{figures/}{./}}
\hypersetup{pdftitle={基于相容区域与覆盖认证的无线电干扰源定位清除},pdfauthor={},pdfsubject={},pdfkeywords={有界误差,半平面交,最坏直径,方向覆盖}}
\setlength{\emergencystretch}{2em}
\setlength{\tabcolsep}{5pt}
\labelformat{figure}{#1}
\labelformat{table}{#1}
\labelformat{equation}{#1}
\biaoti{基于相容区域与覆盖认证的\\无线电干扰源快速定位与清除}
\keyword{有界测向误差；半平面交；最坏直径；定向覆盖；滚动路径规划}
\begin{document}
\begin{abstract}
针对机器狗在有限接收距离及有界测向误差下的无线电干扰源定位清除问题，本文以相容位置集合为基础，将交会定位、主动选点与覆盖搜索统一为逐步缩小不确定区域的决策过程，在保证清除完整性的前提下减少移动及检测开销。

针对问题一，将测向扇区转化为半平面交，分别识别空集、无界和退化情形，对有界凸多边形采用旋转卡壳计算直径。通过可实现的等边三角形观测反例说明：以区域直径为直径的圆未必覆盖全域。进一步以最小覆盖圆半径建立20米光学定位的充分必要判据。

针对问题二，利用第一次收信隐含的接收半径下界，推导四圆盘交集安全域；以最坏相容读数下的区域直径为目标，结合解析条带上界、目标圆条件化及读数网格半步增宽进行分层选点。在两个自建场景中，同口径直径上界分别由115.693米和54.024米降至112.329米和15.769米，结果为有限候选中的优选。

针对问题三，建立原点及六个环点的覆盖布局，将搜索、顺路补测和自适应清除联合调度，采用光学网格作为有限兜底。在192例自建留出案例中全部清除并完成确认，平均耗时为276.0秒/源。

针对问题四，将任意定向朝向可见性转化为邻近检测点的凸包条件，利用有限Voronoi圆缺极值检验认证23点布局，结合双阴性截断、历史观测复用与滚动路径优化。在512例新留出案例中全部清除并完成确认，平均耗时484.3秒/源，比同场景旧版减少10.2\%。

四问实现共通过142项本地测试。本文区分数学保证、自建实验和可核对的接口记录；现有证据尚不足以填报两问各三次正式测试成绩。模型在题设条件与足够执行预算下具有完整性依据，但有限搜索和路径启发式不保证全局最优。
\end{abstract}

\section{问题重述}\label{sec:1}
\subsection{问题背景}\label{sub:1.1}
无线电干扰源的精准定位与清除是无线电频谱管理的重要任务。传统作业依赖技术人员携带测向设备徒步排查，在危险环境中存在安全风险，定位效率也受到限制。为实现自动化作业，现拟利用搭载测向机及光学定位、清除装置的机器狗，通过多点方位观测确定干扰源位置。

在完成初步监测和区域性排查后，精准定位仍受到测向误差与有限接收距离的制约，单点观测难以唯一确定干扰源位置。对于数量未知的多目标场景，频繁切换频道、重复检测及不合理移动还会增加作业时间。因此，需要统筹区域搜索、检测点选择与定位清除过程，在保证清除完整性的同时提高整体作业效率。

\subsection{问题提出}\label{sub:1.2}
已知干扰源位于半径为1800米的圆形目标区域内，以圆心为原点，正东、正北分别为横、纵轴正方向，示向度由横轴正向逆时针计量。误差不超过$1^\circ$且同一地点误差固定，有效接收半径在1000至1500米之间、具体值未知。各源频道互异，取自20个频道。围绕自动定位与清除，需解决以下四个问题。

\textbf{问题一：}根据若干检测点坐标及其对同一干扰源的示向度，设计交会定位区域直径的计算算法，并判断以该直径为直径的圆能否覆盖整个区域。

\textbf{问题二：}对于全向干扰源，在已知一个检测点及其示向度的条件下，制定第二检测点的选择策略，并给出能够获得较好定位效果的候选区域。

\textbf{问题三：}区域内有10至16个全向源，具体数量未知。机器狗由原点出发、初始频道为1，制定自动搜索、定位与清除策略，在确保全部清除的前提下尽量缩短总任务时间，并通过模拟器演练评价清除比例和平均定位清除时间。

\textbf{问题四：}区域内同时存在全向与定向源，总数仍为10至16个，定向源数量与方向均未知。定向源只在其发射方向两侧各$90^\circ$内辐射，其余条件与问题三相同。制定适应混合场景的检测清除策略，并通过模拟器检验。

问题三、四均要求在充分演练后分别完成三次正式测试，按题设记录案例编码、清除数量、平均定位清除时间及程序运行时间，并将原名加密日志纳入支撑材料。

\section{问题分析}\label{sec:2}
\subsection{问题一}\label{sub:2.1}
考虑测向误差后，单次观测将源位置约束在示向线两侧各$1^\circ$的角扇区内，多次观测的交集即为交会定位区域。可将各扇区转化为两个半平面约束，先判定交集状态，再利用凸性把直径转化为最远顶点对距离。覆盖性须独立判断：最大跨度不能直接确定覆盖全域所需的圆半径，应结合反例与最小覆盖圆作答。

\subsection{问题二}\label{sub:2.2}
第二点选择需兼顾再次收信与交会几何。两次测向方向接近共线会放大位置歧义，而任意横向移动又可能超出未知接收范围。因此，先利用第一次收信信息构造保守安全域，再结合目标圆约束，以第二次观测后区域直径的最坏上界评价选点。用分层搜索给出优选点及近优候选区域。

\subsection{问题三}\label{sub:2.3}
源数量未知，单纯追逐已发现目标无法证明任务完成，必须保留覆盖全域的搜索依据。全向源的接收半径下界可将发现过程转化为有限圆盘覆盖。由于移动与逐频道检测均有成本，应将发现新源和清除已知源联合调度，并利用途中停点补充观测。对于几何交会效果较差的区域，采用有限光学覆盖保证已知源最终能够被清除。

\subsection{问题四}\label{sub:2.4}
定向源失收可能源于距离过远，也可能源于处在发射背面，故不能直接照搬全向覆盖或将阴性点周围圆盘删除。应构造对所有朝向均有效的检测点集，并为连续目标域提供覆盖证明。在定位阶段，仅在接收距离已获保证时，结合阳性点与成对阴性点提取几何约束；随后通过历史观测复用与滚动路径改进降低额外成本。

\section{模型的假设}\label{sec:3}
在题设基础上，采用以下假设与计算约定。

\textbf{（1）平面静态场景。}忽略高度差，在二维平面内建模；一次任务期间源位置、频道、发射方向及接收半径保持不变。

\textbf{（2）点位与运动准确。}机器狗能够准确获知并到达指定检测点，路径可无障碍直达；忽略自身定位误差、转向和加减速额外耗时，直线速度为5米/秒。

\textbf{（3）固定有界误差。}不设测向误差概率分布或独立性，不利用同点重复观测平均缩小误差界。前两问按$1^\circ$建模；第三、四问将两位小数读数的舍入误差并入，取$1.005^\circ$。

\textbf{（4）近源与清除。}信号覆盖内距源不超过5米时采用近源反馈；光学定位与清除仅受20米距离条件限制。检测、换频分别耗时5秒、1秒；光学检测耗时3秒，成功后激光清除另计2秒。

\clearpage
\section{符号说明}\label{sec:4}
\begin{table}[H]
\centering
\begin{tabularx}{\textwidth}{cXc}
\toprule
符号 & 含义 & 单位\\
\midrule
$\Omega$ & 半径为1800米的目标圆域 & --\\
$G,g$ & 干扰源的真实位置（仅用于数学分析与事后评价） & 米\\
$S_i,p_i$ & 第$i$个检测点位置 & 米\\
$\theta_i,\varepsilon$ & 示向度及采用的测向误差界 & 度\\
$R$ & 未知有效接收半径，$1000\le R\le1500$ & 米\\
$P,U$ & 与已有观测相容的位置集合或其保守外包 & --\\
$D(P),R_*(P)$ & 区域直径、最小覆盖圆半径 & 米\\
$q=(a,b)$ & 第二检测点在第一示向局部坐标中的位置 & 米\\
$J(q),A(q),B(q),C(q)$ & 最坏直径目标及其解析、几何、组合上界 & 米\\
$V_3,V_4$ & 全向与混合场景的有限覆盖检测点集 & --\\
$N_c,N$ & 已清除数、真实总源数（$N$不输入策略） & 个\\
$T,\overline T$ & 虚拟总任务时间、每源平均定位清除时间 & 秒\\
\bottomrule
\end{tabularx}
\end{table}
其他局部符号在第一次使用时说明。$B(p,r)$表示以$p$为圆心、$r$为半径的圆盘，与第二问上界函数$B(q)$按参数形式区分。

\section{模型的建立与求解}\label{sec:5}
本文以所有与题设及反馈相容的位置构成定位结果，将新增观测转化为集合约束。误差扇区的半平面表达与有界方位定位研究相一致\cite{gholami2015}。第一问处理几何评价，第二问处理单次主动选点，第三、四问进一步将集合更新与有限覆盖和动作调度结合。

\subsection{问题一：半平面交会与定位区域直径}
\label{sub:model_p1}
第一问的算法流程如图\ref{fig:p1_algorithm}所示。先将观测转化为半平面约束并判定区域状态，再对有界交集计算直径，最后检验直径圆的覆盖性与20米光学定位条件。
\begin{figure}[htbp]
 \centering
 \includegraphics[width=.96\textwidth]{p1_algorithm_flowchart.pdf}
 \caption{问题一的半平面交会、区域直径与覆盖判定流程}
 \label{fig:p1_algorithm}
\end{figure}

\subsubsection{示向约束及区域状态判定}
设第$i$个检测点为$S_i$，示向度为$\theta_i$，误差界为$\varepsilon=1^\circ$，定义边界方向
\begin{equation}
 u_i^{\pm}=\bigl(\cos(\theta_i\pm\varepsilon),\sin(\theta_i\pm\varepsilon)\bigr)
\end{equation}
记二维叉积为$[a,b]=a_xb_y-a_yb_x$，则该次观测对应的角扇区为
\begin{equation}
 W_i=\{X:[u_i^-,X-S_i]\ge0,\ [u_i^+,X-S_i]\le0\},\qquad
 P=\bigcap_{i=1}^{n}W_i=\{X:AX\le b\}
 \label{eq:angular_intersection}
\end{equation}
上述向量形式直接约束源位置相对于两条扇区边界的方位，无需对跨越$0^\circ$的角区间另行拆分，使全部观测能够统一表示为线性不等式组。

为避免将空集或无界区域误判为有限多边形，首先枚举不平行边界线的交点，保留满足全部约束者；若未找到可行交点，再检查原点及原点在各边界上的正交投影。该判定是完备的：非空闭凸集距原点最近的点，必为原点、单条有效边界上的垂足或两条独立边界的交点。其次，对非空交集检验
\begin{equation}
 \exists d\ne0,\quad Ad\le0
 \label{eq:recession}
\end{equation}
若成立，则从任一可行点沿$d$无限延伸仍保持可行，区域无界；否则区域有界。二维实现仅需检查各非零约束法向量对应的两条切向方向。对有界交集取凸包，根据顶点数区分单点、线段和多边形。

\subsubsection{区域直径的计算}
设有界区域的顶点为$v_1,\ldots,v_k$。任意$p,q\in P$均可写成顶点的凸组合，从而
\begin{equation}
 \|p-q\|\le\sum_{i,j}\alpha_i\beta_j\|v_i-v_j\|
 \le\max_{i,j}\|v_i-v_j\|,
 \qquad D(P)=\max_{i,j}\|v_i-v_j\|
 \label{eq:vertex_diameter}
\end{equation}
因此只需求最远顶点对。本文采用旋转卡壳方法\cite{toussaint}，在逆时针凸包上依次更新对踵支撑点；遇到平行支撑边时同时检查两端点。直径阶段需要$O(k)$次算术操作，边界交点枚举及约束检验仍为$O(m^3)$，其中$m$为约束数，故不能将完整算法视为线性复杂度。

以$S_1=(0,100)$、$S_2=(100,0)$，示向度分别为$0^\circ$、$90^\circ$的双站观测为自建算例。如图\ref{fig:p1_orthogonal}所示，两扇区的交集为四边形，最远顶点对为$V_2,V_4$，计算得$D=4.9385$米。左图展示测向关系，右图放大交会区域，说明由扇区求交到顶点对计算的几何过程。
\begin{figure}[H]
 \centering
 \includegraphics[width=\linewidth,trim=10 35 10 40,clip]{p1_orthogonal_bearings.pdf}
 \caption{双站正交测向及交会区域（自建算例）。左图为检测点与示向扇区，右图为交会区域及其直径、最小覆盖圆的放大图。}
 \label{fig:p1_orthogonal}
\end{figure}

\subsubsection{直径圆的覆盖判定}
区域直径描述候选位置的最大间距，能否据此构造覆盖全域的圆，仍须另行检验。若$a,b$为最远点对，任何半径$D/2$且覆盖$a,b$的圆，其圆心只能是$c_D=(a+b)/2$。由圆盘的凸性，存在直径为$D$的覆盖圆当且仅当
\begin{equation}
 \max_j\|v_j-c_D\|\le D/2
 \label{eq:diameter_circle}
\end{equation}
图\ref{fig:p1_orthogonal}中的双站算例满足该条件，其最小覆盖圆半径为$2.4693$米，恰为$D/2$；这一关系并不对一般交会区域成立。

为给出与合法观测相容的反例，构造等边三角形$A=(0,0)$、$B=(40,0)$、$C=(20,20\sqrt3)$，取三个检测点
\begin{equation}
 S_1=(-1000,0),\quad S_2=(540,-500\sqrt3),\quad
 S_3=(520,520\sqrt3),
\end{equation}
示向度依次为$1^\circ,121^\circ,241^\circ$。三条下边界对应三角形的有向边，另一侧最大偏角为$\arctan(20\sqrt3/1020)<2^\circ$，故三扇区交集恰为该三角形。真实源可取其重心，接收半径取1200米，与题设相容。

如图\ref{fig:p1_triangle}所示，该区域的直径$D=40$米。以$AC$为直径的圆半径为20米，而顶点$B$到其圆心的距离为$20\sqrt3>20$米，故该圆不能覆盖全域。由最远点对对圆心的唯一约束，任何直径为40米的圆均不能覆盖该三角形，故以区域直径为直径的圆不一定覆盖整个定位区域。
\begin{figure}[H]
 \centering
 \includegraphics[width=.78\linewidth]{p12_geometry.pdf}
 \caption{直径圆不能覆盖定位区域的反例（自建算例，边长40米）。以$AC$为直径的圆未覆盖顶点$B$；最小覆盖圆半径为$40/\sqrt3\approx23.094$米，大于20米。}
 \label{fig:p1_triangle}
\end{figure}

\begin{samepage}
为建立正确的覆盖判据，进一步计算最小覆盖圆，其边界由至多三个点确定\cite{welzl}。单点情形取该点为圆心、半径为0；其余情形枚举两点直径圆与三点外接圆，选择覆盖全部顶点的最小者；该确定性实现最坏需要$O(k^4)$次算术操作。记其半径为$R_*$，平面上的Jung不等式\cite{jung}给出
\begin{equation}
\frac{D}{2}\le R_*\le\frac{D}{\sqrt3}
 \label{eq:jung}
\end{equation}
\end{samepage}
下界由最远点对立即得到；上界亦可由支撑点说明：两点支撑时$R_*=D/2$，三点支撑时圆心在支撑三角形内，其最大角$\gamma\in[60^\circ,90^\circ]$，由正弦定理$R_*\le D/(2\sin\gamma)\le D/\sqrt3$。
因此，一次20米光学定位能够保证覆盖全部候选源的准确判据是$R_*\le20$；仅知直径时，$D\le20\sqrt3$是充分条件，而$D\le40$仅为必要条件。图\ref{fig:p1_triangle}中$R_*=40/\sqrt3>20$米，也直接说明$D=40$米仍不足以保证一次光学定位覆盖全域。

\begin{table}[htbp]
 \centering
 \caption{第一问自建算例的几何结果}
 \label{tab:p1_geometry}
 \begin{tabular}{lrrc}
 \hline
 算例 & $D$/米 & $R_*$/米 & 直径为$D$的圆可覆盖\\
 \hline
 双站正交交会 & 4.9385 & 2.4693 & 是\\
 边长20米的三角形 & 20.0000 & 11.5470 & 否\\
 边长40米的三角形 & 40.0000 & 23.0940 & 否\\
 较宽基线交会 & 31.4797 & 15.7398 & 是\\
 \hline
 \end{tabular}
\end{table}
平行发散观测与不相容观测分别得到无界区域和空集。
\FloatBarrier

\subsection{问题二：保证收信的第二检测点选择}
\label{sub:model_p2}
\subsubsection{第一次检测信息与最坏直径目标}
题设仅给出测向误差界，本文以与观测及题设约束相容的全部源位置表示定位结果，不额外指定其概率分布。
取第一次检测点$S_1$，沿第一示向方向建立局部正交坐标，令$u=(\cos\theta,\sin\theta)$、$n=(-\sin\theta,\cos\theta)$，第二点写为$S(q)=S_1+a u+b n$，其中$q=(a,b)$。第一次正常读数对应的候选源集合为
\begin{equation}
 U_1=\{G:\|G\|\le1800,\ 5<\|G-S_1\|\le1500,
 \ |\operatorname{wrap}(\arg(G-S_1)-\theta)|\le\varepsilon\}
\end{equation}
这里$\operatorname{wrap}$表示取环形角差。固定第一次检测信息后，记$\Psi(q)$为全部相容的第二次正常读数，收到读数$\psi$后保留下来的源位置集合为
\begin{equation}
 \begin{aligned}
 U(q,\psi)=U_1\cap\{G:\;&5<\|G-S(q)\|\le1500,\\
 &|\operatorname{wrap}(\arg(G-S(q))-\psi)|\le\varepsilon\}
 \end{aligned}
\end{equation}
两次正常读数与同一未知半径相容，当且仅当$\max\{1000,\|G-S_1\|,\|G-S(q)\|\}\le1500$，故上述距离条件已包含这一信息。采用Tokekar和Isler的最坏扇区交集直径准则\cite{tokekar}，并将其条件化于已知第一次检测信息，建立
\begin{equation}
 J(q)=\sup_{\psi\in\Psi(q)}\operatorname{diam}U(q,\psi),\qquad
 \min_{q\in\mathcal F}J(q)
 \label{eq:p2_objective}
\end{equation}
该指标与第一问一致，衡量最不利合法读数下的剩余位置歧义。
直径刻画候选位置间的最大分离，可反映面积较小但形状细长的区域仍具有的位置歧义。候选点须对全部相容源位置及接收半径保证收信，这构成鲁棒可行性约束\cite{bertsimas2011}；在该约束下再比较最坏定位损失。以下分别推导安全候选域及两类评分上界。

\subsubsection{四圆盘安全候选域}
设第一距离为$r$、未知接收半径为$R\in[1000,1500]$。第一次收信还意味着$R\ge r$，故保证第二点对所有相容半径仍能收信的条件为
\begin{equation}
 \|S(q)-G\|\le\max(1000,r)
\end{equation}
令$u_\pm=(\cos\varepsilon,\pm\sin\varepsilon)$，暂将第一次检测信息放松为完整环扇形$5<r\le1500$、$|\phi|\le\varepsilon$，则连续安全区域可写为
\begin{equation}
 \mathcal C_{\rm safe}
 =\bigcap_{d\in\{5,1000\},\ s\in\{-,+\}}
 \{q:\|q-du_s\|\le1000\}
 \label{eq:four_disks}
\end{equation}
证明如下：$d=1000$的两个圆盘给出$q\cdot u_\pm\ge\|q\|^2/2000\ge0$，因此$q\cdot u_\phi$在允许角度内的最小值位于角端点。对$5\le r\le1000$，距离平方关于$r$为凸函数，只需检验两个径向端点；对$1000\le r\le1500$，要求化为$\|q\|^2-2r q\cdot u_\phi\le0$，最严格处为$r=1000$。由此得到四个圆盘。目标圆进一步限制源的位置，故该区域对实际问题仍安全，但不一定最大。

图\ref{fig:p2_four_disks}展示式\eqref{eq:four_disks}中的四个圆盘。蓝色、金色分别对应径向端点5米和1000米，实线、虚线区分正负角端点；灰色交集为$\mathcal C_{\rm safe}$。近端两圆心相距约0.175米，在图示尺度下近乎重合。
\begin{figure}[htbp]
 \centering
 \includegraphics[width=.78\textwidth]{p2_02_four_disks.pdf}
 \caption{四圆盘交集构成的保证收信区域$\mathcal C_{\rm safe}$}
 \label{fig:p2_four_disks}
\end{figure}


若直接横向移动，取真实源$(1000,0)$、$R=1000$，任意第二点$(0,b)$、$b\ne0$均会失收，说明改善交会夹角必须服从收信约束。为统一比较第二次正常测向，本文采用
\begin{equation}
 \mathcal F=\mathcal C_{\rm safe}\cap
 \{q:\|q\|\le L,\ |b|\cos\varepsilon-a\sin\varepsilon>5\}
 \label{eq:p2_feasible}
\end{equation}
最后一项保守排除了第二点进入5米近区的情况；实际遇到近区反馈可直接清除。默认$L=1100$米，四圆盘本身已蕴含$\|q\|\le1005$米。1800米目标圆仅约束干扰源，第二检测点可以位于圆外。

\subsubsection{评分方法一：解析条带上界}
参照Tokekar和Isler的几何外包思路\cite{tokekar}，本文将两次测向扇区分别放松为条带，以交集的长对角线构造解析评分$A(q)$。图\ref{fig:p2_strip}中两条带的半宽分别为
\begin{equation}
 \begin{aligned}
 w_1&=1500\sin\varepsilon,\qquad w_2=r_{2,\max}\sin\varepsilon,\\
 r_{2,\max}&=\max_{r\in\{5,1500\},\ s\in\{-,+\}}\|q-r u_s\|
 \end{aligned}
 \label{eq:p2_strip_widths}
\end{equation}
\begin{figure}[H]
 \centering
 \includegraphics[width=.78\textwidth]{p2_08_strip_intersection.pdf}
 \caption{评分一的条带交集与解析直径：夹角取下界$\beta_0$，双向箭头沿长对角线标示$A(q)$}
 \label{fig:p2_strip}
\end{figure}
图中$O$为条带中心线交点，$\xi,\eta$为平移后的示意坐标。为使直径界对全部相容的第二读数成立，令
\begin{equation}
 \begin{aligned}
 h_0&=|b|-w_1,\qquad k_0=\max\{|a-5\cos\varepsilon|,|a-1500|\},\\
 \beta_0&=\arctan2(h_0,k_0)-\varepsilon
 \end{aligned}
 \label{eq:p2_crossing_lower}
\end{equation}
其中$h_0$给出第二检测点到第一条带的垂直间隔，$k_0$界定其与候选源的水平间隔。当$h_0>0$且$\beta_0>0$时，两条示向中心直线的锐交会角不小于$\beta_0$。

固定半宽时，条带交集为平行四边形，其直径随锐交会角增大而减小。取夹角下界$\beta_0$，由图中长对角线$V_1V_3$及余弦定理得
\begin{equation}
 J(q)\le A(q)=
 \frac{2\sqrt{w_1^2+w_2^2+2w_1w_2\cos\beta_0}}{\sin\beta_0}
 \label{eq:analytic_bound}
\end{equation}
若$h_0\le0$或$\beta_0\le0$，则置$A(q)=+\infty$。

\subsubsection{评分方法二：几何外包上界}
沿用Tokekar和Isler以相容测向扇区交集的最坏直径评价定位质量的方法\cite{tokekar}，本文进一步保留目标圆和距离约束，将第二读数划分为单元，分别构造几何外包，再取其直径的最大值作为评分$B(q)$。

记第一扇区、第一条带和第一距离下界的有效投影约束分别为$W_1,T_1,H_1$；$D_0$为目标圆盘，$D_1,D_2$为第一、第二距离上界圆盘，其360边外切多边形记为$\widehat D_i$。将这些约束求交，构造与第二读数无关的凸多边形外包$Q(q)$，其集合关系见图\ref{fig:p2_reading_cells}。这里对5米近区采用前向投影下界放松，保留真实相容位置。

令读数步长为$h=360^\circ/N$，半步为$\delta=h/2$，第$k$个格中心为$\psi_k$，以与$\psi_k$的环形角差不超过$\delta$的角度组成读数单元$I_k$。将$Q(q)$与该单元的增宽测向约束求交，得到覆盖单元内全部相容读数的外包$P_k(q)$，具体增宽方式见式\eqref{eq:half_step}。图\ref{fig:p2_reading_cells}以自建原点场景的推荐点$q_*$为例，展示这一构造。
\begin{figure}[H]
 \centering
 \includegraphics[width=.96\textwidth]{p2_09_reading_cells.pdf}
 \caption{同一候选点下的读数单元、几何外包及圆边界近似}
 \label{fig:p2_reading_cells}
\end{figure}
上方主图的灰色区域为$Q(q_*)$，三条射线对应格中心$231.0^\circ$、$276.0^\circ$和$318.0^\circ$。图\ref{fig:p2_reading_cells}(a)--(c)分别放大同色单元外包，黑色线段标示直径$D_{k_i}$，其值依次为46.60米、39.28米和110.64米。图(a)、(b)中的外包均位于三个真实圆盘内；图(c)则受到第一距离圆盘外切边的裁切，圆圈标出该处边界。

图\ref{fig:p2_reading_cells}(d)、(e)分别展示目标圆与第二距离圆的外切边界构造，并非图(a)、(b)的局部放大；图\ref{fig:p2_reading_cells}(f)对应图(c)圈出处的实际裁边。图(f)中约1.884毫米的差异仅表示该单元局部的圆边外包增量，不是定位误差或全部顶点的误差上限。最下排采用切向米、法向毫米的非等比例坐标，以分辨圆弧与外切边的细微差别。

由图(a)--(c)可见，检测点固定时，不同读数仍会产生形状和大小不同的外包区域。因此须评价全部单元的最坏情况：
\begin{equation}
 D_k(q)=\operatorname{diam}P_k(q),\qquad
 J(q)\le B(q)=\max_k D_k(q)
 \label{eq:geometric_bound}
\end{equation}
空外包约定$D_k=0$；三格示例不能代替完整的$B(q)$，$D_k$不是真实定位误差。

上述上界成立的关键是覆盖连续读数。任意读数距最近格中心至多$\delta$，故须同步增宽所有随读数转动的约束：
\begin{equation}
 \begin{aligned}
 \text{扇区半角}&:\ \varepsilon+\delta,\qquad
 \text{前向投影下界}:\ 5\cos(\varepsilon+\delta),\\
 \text{条带半宽}&:\ R_2\sin\varepsilon+2R_{\rm out}\sin(\delta/2),\\
 R_2&=\min(1500,r_{2,\max}),\qquad
 R_{\rm out}=\max_{v\in\operatorname{vert}Q(q)}\|v-S(q)\|
 \end{aligned}
 \label{eq:half_step}
\end{equation}
条带增量来自旋转前后单位法向量的差范数至多$2\sin(\delta/2)$。投影下界也须随角度放松：若原相对角为$t\in[-\varepsilon,\varepsilon]$且$r\cos t\ge5\cos\varepsilon$，读数旋转$|\Delta|\le\delta$后的投影至少为$5\cos(\varepsilon+|\Delta|)\ge5\cos(\varepsilon+\delta)$。该推导同时覆盖原始外包中因近区放松引入的点。于是，对任意相容读数$\psi\in I_k$，均有$U(q,\psi)\subseteq P_k(q)$；结合直径关于集合包含关系的单调性，即得式\eqref{eq:geometric_bound}。

评分单元与第二读数后的区域共用上述几何构造。默认$h=0.1^\circ$，图中评分单元的扇区半角为$1.05^\circ$；获得实际读数后仍使用$1^\circ$。

\subsubsection{评分组合与分层选点}
两种方法分别给出同一最坏直径目标的上界，故取二者较小值作为最终评分：
\begin{equation}
 C(q)=\min\{A(q),B(q)\}\ge J(q)
 \label{eq:combined_bound}
\end{equation}
较小的$C$表示较紧的最坏保证，未必意味着未知精确值$J$也按同样顺序排列。

在候选搜索层面，借鉴DIRECT算法\cite{jones1993}与多层坐标搜索（MCS）\cite{huyer1999}兼顾全域探索和局部细化的思想，本文采用极坐标粗网格搜索，并围绕当前优选点及其镜像方位逐级细化。

在式\eqref{eq:p2_feasible}限定的候选域内，依次采用100米/$5^\circ$、20米/$1^\circ$及5米/$0.25^\circ$的移动距离与方位步长搜索，并加入所搜索方位的安全边界点。三层的读数评分步长分别为$1^\circ$、$0.5^\circ$和$0.25^\circ$；各层保留既有候选，最后以$0.1^\circ$统一复核。图\ref{fig:p2_progressive_zoom}以第一检测点为$(0,0)$、第一示向为$0^\circ$的自建算例展示这一过程。

\begin{figure}[htbp]
 \centering
 \includegraphics[width=\textwidth]{p2_14a_progressive_zoom.pdf}
 \caption{原点自建算例的三层候选搜索与逐级放大}
 \label{fig:p2_progressive_zoom}
\end{figure}

图\ref{fig:p2_progressive_zoom}中蓝、金、绿点分别表示各层首次访问的候选，灰色空心点为保留的旧候选，星号为本层评分精度下的优选点。灰底为安全候选域，虚线框及引线指示下一子图的放大范围。各子图采用原始局部坐标，横纵等比例，仅展示上半域的全部新增点；下半域实际访问点与其逐点镜像对应。因此，逐级收紧的是新增加密网格，累计候选集仍随层数增加。

三层在完整候选域中分别新增320、136和92个点，累计访问548点；图中对应显示160、68和46个新增点。由于各层读数评分精度不同，第三层星号不代表最终推荐点。本例统一复核后选回第二层首次访问的$(842.144,546.895)$米候选，其评分上界为112.3293米，与表\ref{tab:p2_selection}一致。

统一复核时，候选点按精度优先、距离次优的字典序筛选\cite{mavrotas2009}：先比较$C(q)$，仅评分平局时优先较短移动；实现按$10^{-7}$米舍入判定数值平局。记推荐点为$q_*$，给定可调容差$\eta$，近优候选区域定义为
\begin{equation}
 \mathcal C_{\rm good}=\{q\in\mathcal F:C(q)\le(1+\eta)C(q_*)\},\qquad \eta=10\%
\end{equation}
数值输出给出该连续区域中已访问的代表点。
另以$\mathcal V$表示已访问候选集，提供
\begin{equation}
 q_\eta\in\arg\min_{q\in\mathcal V\cap\mathcal C_{\rm good}}\|q\|
 \label{eq:p2_epsilon_selection}
\end{equation}
这采用$\epsilon$-约束思想\cite{mavrotas2009}，将定位要求写成阈值，再缩短路程。由于本次移动与检测时间为$\|q\|/5+5$，缩短路程等价于减少该项时间。

计算时先求$A$，依次扫描$B$的单元直径。已扫最大值$b_t$单调不减，一旦$b_t\ge A$，便已确定$C=A$，可提前停止。

安全候选生成、组合评分及分层选点的完整流程见图\ref{fig:p2_algorithm}。各层保留已访问点，最终统一复核，输出推荐点及近优代表点。
\begin{figure}[htbp]
 \centering
 \includegraphics[width=.98\textwidth]{p2_algorithm_flowchart.pdf}
 \caption{问题二的安全候选生成、组合评分与分层选点流程}
 \label{fig:p2_algorithm}
\end{figure}


\begin{table}[htbp]
 \centering
 \caption{第二问自建场景的同口径比较（均保证收信）}
 \label{tab:p2_selection}
 \begin{tabular}{llrrr}
 \hline
 第一检测点/米 & 策略 & 局部第二点/米 & $C$/米 & 移动加检测/秒\\
 \hline
 $(0,0)$ & 解析基线 & $(750,660)$ & 115.6930 & 204.810\\
 $(0,0)$ & 本文选点 & $(842.144,546.895)$ & 112.3293 & 205.828\\
 $(0,0)$ & 10\%内最短点 & $(773.511,541.618)$ & 122.8653 & 193.857\\
 $(1500,0)$ & 解析基线 & $(750,660)$ & 54.0236 & 204.810\\
 $(1500,0)$ & 本文选点 & $(220.942,195.473)$ & 15.7688 & 64.000\\
 \hline
 \end{tabular}
\end{table}
表\ref{tab:p2_selection}的第一示向度均为$0^\circ$，圆外包取360边、读数步长取$0.1^\circ$，时间仅计本次移动及5秒检测。相对于解析基线点的同口径评分，原点和边界场景的上界分别下降2.91\%和70.81\%；这表示最坏直径保证变紧，不等于真实定位误差同比下降。在原点算例中，对61个第一距离与21个第一方位组成的1281组源位置逐一验证，推荐检测点均能接收到信号，与四圆盘安全域的理论判定一致。

完整扫描逐一计算全部读数单元的直径；提前停止则在已扫最大值达到解析上界$A$时终止余下扫描。默认原点流程访问548个候选，两种方式得到相同的逐层优选点、最终评分及排序。扫描量由3\,205\,440降至3\,201\,498，节省0.123\%；两种方式的完整流程耗时分别为14.821秒和14.868秒，实测运行时间基本相当。固定本文原点选点，将精度提高至720边、$0.05^\circ$后，上界由112.3293降至111.7932米，反映外包与读数网格仍有保守性。
% 正文片段：置于“模型的建立与求解”中；末尾检验片段可移至检验章节。
\subsection{问题三：全向干扰源的覆盖搜索与自适应清除}
\label{sub:model-p3}

\subsubsection{时间目标与相容区域}
将任务分为发现、定位和清除三个相互衔接的过程。设总路程为$L$，检测次数、切换频道次数、失败清除次数及成功清除次数分别为$N_m,N_s,N_f,N_c$，按题设计时规则得
\begin{equation}
 T=\frac{L}{5}+5N_m+N_s+3N_f+5N_c,
 \qquad \overline T=\frac{T}{N_c},\quad N_c>0
 \label{eq:search-time}
\end{equation}
优化目标是在全部清除并取得完成依据的条件下减小$T$。清除指令不改变测向机频道，故频道切换须按实际相邻检测操作计算；最后一次清除之后的覆盖确认时间仍计入总时间。真实源数仅在自建实验结束后用于评价，不向策略提供。

对每个已发现频道$j$，维护可能位置集合$P_j$。初始集合为目标圆$\Omega$；在$p_i$得到示向度$\theta_i$后更新为
\begin{equation}
 P_j\leftarrow P_j\cap S(p_i,\theta_i,\varepsilon)
          \cap B(p_i,1500),\qquad \varepsilon=1.005^\circ,
 \label{eq:search-region}
\end{equation}
其中$S$为闭角扇区，$B(p,r)$为闭圆盘。误差界包含题设$1^\circ$及接口两位小数舍入的半个末位单位。外接正多边形构造在实数运算下只放宽候选集合；浮点实现配合向外容差，并通过边界回归检验。同一地点的重复读数不作平均。无信号反馈不直接缩小此处的候选多边形，以保持观测不确定时的保守性。

\subsubsection{七点覆盖与未知源发现}
设搜索点集为
\begin{equation}
 V_3=\{(0,0)\}\cup\left\{1150
 \left(\cos\frac{k\pi}{3},\sin\frac{k\pi}{3}\right):k=0,\ldots,5\right\}
 \label{eq:p3-points}
\end{equation}
在每个搜索点检测尚未发现、尚未清除的频道，优先检测当前频道以减少切换。原点与六个环点的内部Voronoi顶点距最近搜索点不超过$1150/\sqrt3$；目标圆边界上的最不利位置位于相邻环点夹角的平分线上。因此覆盖半径满足
\begin{equation}
 r_{\rm cov}\le\max\left\{\frac{1150}{\sqrt3},
 \sqrt{1800^2+1150^2-2\cdot1800\cdot1150\cos30^\circ}\right\}<1000
 \label{eq:p3-cover}
\end{equation}
由于实际接收半径不小于1000米，每个全向源至少在其中一点被发现。由此将连续区域的搜索转化为有限搜索点的访问，避免仅凭长时间无信号认定搜索结束。

\subsubsection{搜索与定位的联合调度}
单纯先扫完全部搜索点再逐源清除，容易在已走过的区域之间反复移动。本文每次同时比较未知区域的搜索任务与已知源的清除任务。对已知源$j$，用候选多边形的质心$\widehat g_j$及外包圆半径$r_j$构造排序代价
\begin{equation}
 C_j=\|p-\widehat g_j\|+0.5r_j
 \label{eq:p3-cost}
\end{equation}
选择$C_j$最小的源，并与最近待访搜索点的距离比较；当$C_j$不超过该距离加180米时优先定位，否则继续覆盖搜索。半径项抑制对位置尚不明确的源过早追逐。该代价仅用于行动排序，不作为误差界或最优性结论。

定位时沿最近示向线的法向，在候选中心附近添加45至140米的横向位移，以接近目标并形成交会角；已有多次读数且区域较小时，直接接近候选中心检测。每次搜索或清除停点还检查其他已知源：若其候选区外包半径不小于35米、预计距离不超过1450米且预计视差至少为$12^\circ$，则顺路补测。这样利用已付出的移动成本换取额外方位约束，在检测次数和折返路程之间进行权衡。

\subsubsection{可靠清除与有限兜底}
对候选多边形构造外包圆，并重新核对全部顶点。若外包半径$r_j\le19.99$米，则在圆心清除可覆盖真源。这里使用的是覆盖全域的圆半径，不能以区域直径的一半替代。若$r_j\le55$米，可先在质心试清一次以降低平均耗时；此时成功不作事先保证，失败后仍保留原区域。

主动检测最多进行5次。若仍未完成，则沿候选区最长弦方向铺设边长$h=26$米的方格，保留所有与候选区相交的方格中心，依次进行光学清除。任一格内点到中心的距离满足
\begin{equation}
 d\le\frac{h}{\sqrt2}=13\sqrt2\approx18.385<20\ \mathrm{m}
 \label{eq:optical-cover}
\end{equation}
因此这些有限清除点覆盖整个候选区；边缘方格即使中心位于区域外也必须保留。以最近邻顺序访问仅改变耗时，不改变覆盖性质。出现距离过近反馈时，则直接在当前点清除。

在题设误差与设备条件成立、接口执行正确且时间预算足够时，式\eqref{eq:p3-cover}保证全部源被发现，式\eqref{eq:search-region}保证真源始终保留，式\eqref{eq:optical-cover}保证已发现源经有限操作被清除。故终止条件设为：覆盖搜索全部完成且已发现源全部清除，或已实际清除16个不同频道的源。达到10个源不能据此终止。

整体算法可归纳为以下步骤：初始化搜索点、未知频道集合和当前位置；比较最近搜索点与已知源的定位代价并执行较优行动；将有效示向度转化为候选区约束，近源反馈则立即清除；在停点补测具有足够视差的其他已知源；依据区域大小选择直接清除、继续交会或光学兜底；最后检查两类终止条件，未满足时返回行动选择步骤。整个过程只使用已经发生的观测，定位失败不更新清除集合。

\subsection{问题四：混合定向源的认证覆盖与观测复用}
\label{sub:model-p4}

\subsubsection{任意朝向的覆盖条件}
定向源的有效区域为接收圆盘与某个过源位置的闭半平面的交集，无信号还可能由处于发射背面引起。七点全向覆盖不足以保证任意朝向可见，尤其需要关注目标圆边界处朝外发射的源。

\begin{figure}[htbp]
\centering
\includegraphics[width=0.96\linewidth]{search_coverage.pdf}
\caption{第三问全向覆盖与第四问认证检测布局}
\label{fig:search-coverage}
\end{figure}

对源位置$g$，设$N(g)=\{v\in V:\|v-g\|\le1000\}$。任意发射闭半平面中均存在有效检测点的充要条件为
\begin{equation}
 g\in\operatorname{conv}N(g)
 \label{eq:directional-convex}
\end{equation}
若此式成立，写$g=\sum_i\lambda_i v_i$，对任意发射方向单位向量$n$有$\sum_i\lambda_i n^{\mathsf T}(v_i-g)=0$，不可能全部内积为负，故至少一个检测点位于发射半平面内；反之，若$g$在邻点凸包外，由凸集分离可找到不包含任何邻点的发射半平面。

本文采用包含原点、12个外环点和10个非均匀内点的23点布局。固定坐标由几何搜索生成，不依赖实验中的真实源位置。为验证连续目标圆而非有限采样点上的覆盖，对任意两检测点确定的有向直线$\ell$，只取严格在其左侧的检测点$V_\ell^+$，计算
\begin{equation}
 \rho_\ell=\max_{g\in\Omega\cap H_\ell^+}
             \min_{v\in V_\ell^+}\|g-v\|,
 \qquad \max_\ell\rho_\ell<1000,
 \label{eq:directional-certificate}
\end{equation}
其中$H_\ell^+$为左闭半平面，同时要求全体检测点凸包严格包含目标圆。若式\eqref{eq:directional-convex}失效，可由邻点凸包的支撑边构造分隔线，使源位于左侧而全部邻点位于另一侧，违反式\eqref{eq:directional-certificate}；单点或共线等退化情形由全体凸包严格包含条件处理。

计算式\eqref{eq:directional-certificate}时，将圆缺按$V_\ell^+$的Voronoi区域划分\cite{deberg2008}。各单元只需检查多边形顶点、直边与目标圆的交点及圆弧上的距离极大点，即可得到有限极值证书。六位小数坐标的认证结果为：凸包内切半径1800.969米，494个有效有向点对对应的最大接收距离上界984.419米，距1000米下界留有15.581米余量。密集网格仅作实现回归，不作为覆盖证明，也不声称23点为全局最少。

为建立可解释的对照，保留上一版25点布局：原点、半径950米的12个内环点以及半径1870米的12个外环点，内外环错开$15^\circ$。连接相邻环点得到36个三角形。外正十二边形的内切半径为$1870\cos15^\circ>1800$米，而三角形的最大边长约983.598米，故每个源都可写为三个1000米内检测点的凸组合，直接满足式\eqref{eq:directional-convex}。23点方案放宽了“每个三角形均有短边”的构造限制，通过更精细的圆缺极值认证减少搜索点。两种覆盖依据均独立于发射朝向及源分布，检测点允许位于目标圆外。

\subsubsection{成对横向检测与双阴性截断}
阳性观测仍按式\eqref{eq:search-region}更新。设最近阳性点为$p$，示向方向及其法向单位向量为$e,u$，候选区相对$p$的纵向投影满足$e^{\mathsf T}(g-p)\in[a,b]$。取截面$t=a+0.2(b-a)$，在
\begin{equation}
 q_\pm=p+te\pm wu,\qquad
 w\ge\max\{15,\min(140,0.04t)\},\quad w>t\tan\varepsilon
 \label{eq:paired-probe}
\end{equation}
进行成对探测，优先访问距当前位置较近的一侧，出现阳性即更新区域。相较于直接接近区间中点，偏前截面可减少定向源定位过程中的过冲。若两个探点不能同时保证位于整个候选区的1000米范围内，则逐步将截面移向中部，最终仍以显式距离检查决定能否使用阴性反馈。

只有当所有候选顶点到$q_+,q_-$均小于999.999米时，才可排除接收距离不足的原因。在$t>0$时，若候选源满足$e^{\mathsf T}(g-p)\ge t$，线段$[p,g]$与截面的交点必在$[q_-,q_+]$上。阳性点$p$与源$g$均在发射闭半平面内，其连线也在其内；而两个阴性点若都在发射半平面外，它们之间的线段亦在外，产生矛盾。因此双阴性可安全推出
\begin{equation}
 e^{\mathsf T}(g-p)<t
 \label{eq:paired-negative}
\end{equation}
实现时保守地按带余量的闭半平面截断。缺少距离保证、只有一次阴性或探测未完成时均不作此推断。

\subsubsection{历史阴性信息的安全利用}
保存已经访问过的阴性点，得到新阳性观测后重新检查其约束作用。设$A$为历史阳性点，$B,C$为历史阴性点。若两阴性点均可证处于候选源的有效接收半径内，则任何使$[A,g]$与$[B,C]$相交的位置$g$均可排除，理由与双阴性截断相同。

为使长条候选区中的局部信息也能利用，沿最新示向方向将候选区分成12片。对每片及阴性点$Q$，采用以下任一充分条件认证其在接收半径内：其所有顶点到$Q$均小于999.999米；或存在历史阳性点$P$，使所有顶点均满足
\begin{equation}
 2(P-Q)^{\mathsf T}g<\|P\|^2-\|Q\|^2
 \label{eq:relative-reception}
\end{equation}
后一式等价于$\|Q-g\|<\|P-g\|\le R$，利用了历史收信给出的未知半径下界。该判据是关于$g$的线性不等式，顶点成立即可保证整片成立。扣除有依据的禁区后保留所有剩余碎片，再取凸包恢复统一表示；取凸包只会放宽区域。碎片过多时放弃此次可选裁剪，不以效率为由删除未被排除的位置。成对探测最多5轮，之后仍采用式\eqref{eq:optical-cover}的光学覆盖兜底。

\subsubsection{滚动路径规划与全清判定}
将尚未完成的搜索点及已知源候选中心共同作为访问目标，构造从当前位置出发的开放路径。比较最近邻、插入法、正反角度遍历及上一轮未执行路线，经2-opt边交换和连续1至3个访问点的Or-opt搬移改进后，选择较短方案。局部交换遵循旅行商路径改进的基本思想\cite{croes1958}，但这里只执行第一个动作，随后根据观测重新规划，不声称得到全局最短路线。

为兼顾移动与扫频开销，另以固定离散位置、三种半径和八个定向朝向估计未发现频道在后续搜索点上的收信机会，按移动时间加预计扫频时间调整顺序。单频道存在概率取0.65，每次预计检测近似计6秒。该先验仅是排序启发式，不是官方源分布，也不参与相容区域裁剪与全清证明；实际评价仍使用式\eqref{eq:search-time}逐动作累计。

全清终止沿用第三问的两类依据。发现16个不同频道后可以停止发现阶段，但必须将这16个源实际清除才能终止。若提前以源数上界完成任务，应记录“已完成清除证明”，同时保留“未走完全部覆盖点”的事实。有限覆盖与保守定位保证的是模型条件下的完整性，实际接口中断、预算不足等情况必须记为未完成。


\section{模型的检验}\label{sec:6}

\subsection{第三、四问的本地对照与执行记录}
\label{sub:verify-p34}
自建场景包含均匀分布、成簇、边界外向、狭长走廊、近距离双源及最小接收半径等情况，测向误差包括空间相关误差、固定偏置与极端误差。策略仅通过观测接口行动，源真值仅供事后评分。第三问开发、留出、压力集分别为96、192、96例；第四问第二轮采用新的384、512、256例，代码冻结后再评估留出与压力集。表\ref{tab:search-local}中的平均时间先在每例计算$T/N_c$，再对案例取平均。

\begin{table}[htbp]
\centering\small
\caption{当前采用算法在自建样本上的结果}
\label{tab:search-local}
\begin{tabular}{llrrrr}
\toprule
问题 & 集合 & 例数 & 全清且确认 & 平均秒/源 & 阈值通过率\\
\midrule
三 & 开发 &96&96&266.7&66.7\%\\
三 & 留出 &192&192&276.0&61.5\%\\
三 & 压力 &96&96&318.4&45.8\%\\
四 & 第二轮开发 &384&384&483.7&55.5\%\\
四 & 第二轮留出 &512&512&484.3&56.8\%\\
四 & 第二轮压力 &256&256&537.8&43.4\%\\
\bottomrule
\end{tabular}
\end{table}

阈值分别取300、500秒/源，为本地预设比较标准。第三问在相同96例开发样本中，先覆盖后清除的基线、本文策略和直接光学策略分别为502.6、266.7、611.9秒/源。第四问与上一版25点自适应策略逐例比较，新留出集平均时间由538.99降至484.28秒/源，减少10.2\%；压力集由629.14降至537.78秒/源，减少14.5\%。改进同时涉及布局、裁剪和调度，不能将全部收益单独归因于减少两个搜索点。留出集中452例更快、60例更慢，说明提速并非逐例保证；少源情形仍需承担覆盖确认成本，是当前主要瓶颈。

\begin{figure}[htbp]
\centering
\includegraphics[width=0.98\linewidth]{search_comparison.pdf}
\caption{自建样本的平均时间与预设阈值通过率对照}
\label{fig:search-comparison}
\end{figure}

进一步按真实源数进行事后分组，第四问留出集的10源案例平均为624.1秒/源，16源案例为331.8秒/源。一方面，相同覆盖路程在较少源之间分摊时，平均成本自然较高；另一方面，实际清除16个源可利用题设上界提前结束，而少于16个时仍须完成覆盖确认。该差异解释了样本总体全清与速度阈值通过率不高可以同时出现，也说明不能在检测到较少源时过早终止来换取表面上的低耗时。

检验还覆盖测角量化边界、细长候选区域、光学网格边角、定向阴性裁剪的距离前提以及16源提前停止等情形。统一验证记录中的第三、四问分别通过41项和38项测试，并各执行8例困难样本回归；这些检验支持几何实现和状态处理的一致性，不能替代真实接口延迟、连接异常及有限现实预算下的运行检验。

\begin{table}[htbp]
\centering\small
\caption{第三问已匹配的六次平台演练记录}
\label{tab:p3-practice}
\begin{tabular}{lrrr}
\toprule
案例编码 & 已清除数 & 总虚拟时间/秒 & 平均秒/源\\
\midrule
\texttt{WAHP-JBFF-NQX8-CG25}&11&3022.56&274.78\\
\texttt{9B6R-XRTW-CGGF-92KT}&11&3009.29&273.57\\
\texttt{26TK-6AS4-3E9S-RPT5}&13&3631.30&279.33\\
\texttt{WS8S-NUQR-T42H-347E}&13&2795.69&215.05\\
\texttt{PQBH-M69S-H2ZH-UTY7}&15&3499.46&233.30\\
\texttt{WAW2-SF6X-36BJ-A5ZK}&12&3546.01&295.50\\
\bottomrule
\end{tabular}
\end{table}

现有执行目录保存第三问8次、第四问7次接口会话，均确认接受进入与退出请求，且没有未决请求。其中第三问前六次可与留存演练日志的封装头匹配：日志模块标明演练，案例编码如表\ref{tab:p3-practice}，日志创建时间与对应退出响应的服务端时间相差不超过3毫秒。表中清除数量和时间取客户端已接受响应的累计值；日志明文元信息未提供源总数，因而不据此填报平台清除率。其余两次第三问和七次第四问会话尚缺可核对的模块标识与案例编码，仅保留为接口执行记录。上述演练及客户端记录均不替代各问三次正式测试，正式成绩还需核对平台运行时间及原名加密日志。

后续第三轮速度探索仅使用开发样本，尚未启用独立留出及压力集合，不替换表\ref{tab:search-local}的冻结结果，也不将局部开发提升表述为已完成的正式验证。

\subsection{几何一致性与数值敏感性}
第一问分别检验空、无界、单点、线段和一般多边形，核对跨越$0^\circ$的角度、平行边界及等边三角形反例；旋转卡壳结果与顶点对穷举一致。第二问检查四圆盘保证、近源边界、目标圆条件化、第二读数扇区与条带的同步增宽，以及提前停止与完整扫描的评分和选点一致性。四问当前统一验证分别通过36、27、41、38项测试，共142项，并检查异目录启动、原资料完整性及结果可读取性。

误差增大时，真实相容位置集合不会缩小，单个观测的横向不确定性随距离约按$r\sin\varepsilon$增长。因此，接近目标和增加交会夹角通常有利于缩小长条区域，但是否值得补测仍须考虑5秒检测成本。数值敏感性还来自圆盘外包边数与读数步长；第二问的加密复核显示当前上界仍有保守余量。几何检验中的浮点容差可减小边界误删风险，但不等同于区间算术意义上的严格机器证明。

清除比例定义为$N_c/N$，只有在独立获知真实源数时才能计算；每例时间使用$T/N_c$，跨例平均不与$\sum T/\sum N_c$混用。所有本地比较包含最后一次清除之后为确认覆盖完整性付出的时间，防止把提前停止造成的低耗时误认为策略提速。

\subsection{接口执行与正式测试结果的证据边界}
程序按串行HTTP请求执行移动、检测和清除，并为动作配置唯一请求编号。相同动作结果不明时仅以同编号、同内容重试，未决动作没有确认前不发送下一动作；拒绝响应不更新位置、频道或清除集合。进入场景后读取实际剩余程序时间，并以此控制运行预算。清除操作不切换当前测向频道，任务结束后的退出请求不作为查询结束原因的手段。

截至本稿整理时，尚未取得能够完整核对的第三、四问各三次正式成绩，因此不以本地案例或演练数据填充正式表。题设要求的记录结构如下，后续只能依据平台实际结果及对应原名加密日志补入。
\begin{table}[htbp]
\centering\small
\caption{正式测试要求与当前可核验状态}
\begin{tabularx}{\textwidth}{lXXXX}
\toprule
问题及次数 & 测试案例编码 & 清除源数 & 平均定位清除时间 & 程序运行时间\\
\midrule
问题三，三次 & \multicolumn{4}{l}{尚无完整可核对的正式测试记录，未填报成绩}\\
问题四，三次 & \multicolumn{4}{l}{尚无完整可核对的正式测试记录，未填报成绩}\\
\bottomrule
\end{tabularx}
\end{table}
正式测试每问仅三次，成功启动或手工中止均占用机会。论文定稿须逐项核对上述四字段，将日志原名文件纳入支撑材料，并保留平台上传状态。本文已有的模型与本地检验结论不依赖这些尚未补齐的正式成绩。

\section{模型的评价及推广}\label{sec:7}
\subsection{模型的优点}
\textbf{第一，几何含义明确。}以相容集合贯穿四问，区分区域直径、固定点最大误差与最小覆盖圆半径，保证定位评价与20米清除条件一致。无需假定未知误差服从特定分布。

\textbf{第二，完整性与提速依据分明。}四圆盘收信、有限搜索点覆盖和光学网格提供基础保证，滚动路径及顺路补测负责降低平均耗时；后者失效时仍保留基础清除过程。定向阴性信息只有在几何与距离前提均成立时才用于裁剪。

\textbf{第三，可复现且便于检验。}模型、场景生成、覆盖证书及计时规则显式实现；比较使用相同样本、固定地点误差和统一统计口径，并保留更慢案例。参数由开发集选择，独立集合用于后续评价。

\subsection{模型的局限}
第二问连续目标函数采用外包上界和有限选点搜索，最优性仅限已访问候选；第一问当前边界枚举和最小覆盖圆枚举也有较高的最坏复杂度。第三、四问路径规划含多个启发式，不能保证逐例或全局最短时间。23点布局证明了其覆盖有效性，尚未证明点数最少。

自建环境不能代表官方生成分布；少源、边界及较弱交会几何仍可能导致较高确认成本。实现没有显式建模障碍物、自身定位误差、转向动力学或电磁特性随时间变化。接口异常和实际时间预算也可能阻止理论上的有限过程完成。现阶段正式测试记录尚未齐备，无法据此报告正式成绩或对真实场景成功率作统计保证。

\subsection{推广与改进方向}
对于位置误差有界的设备，可先将检测点不确定域纳入测向约束的外包，再重新验证收信与清除条件；若场景存在障碍物，应以可行最短路距离替换欧氏移动成本，并重新认证可访问检测点的覆盖。多机器狗扩展可按频道或区域分配任务、共享相容集合，但须计入通信与冲突成本。

在现有框架内，可通过更紧的圆弧几何、分支界定和连续选点优化降低第二问的外包保守性；对多源任务进一步联合优化搜索点与访问顺序，尤其减少少源情形的覆盖确认成本。任何新增策略均应先冻结参数，再用新的独立样本及平台测试检验，避免将开发集收益直接写成泛化结论。

\section*{AI工具使用声明}
本参赛队在竞赛过程中使用了AI工具，主要用于建模思路讨论、程序实现与调试、实验结果整理、参考文献核查及论文撰写辅助，详细使用情况见支撑材料。

\urlstyle{same}
\begin{thebibliography}{99}
\bibitem{gholami2015} GHOLAMI M R, GEZICI S, WYMEERSCH H, et al. Characterizing the worst-case position error in bearing-only target localization[C]//Workshop on Positioning, Navigation and Communication. 2015.
\bibitem{toussaint} TOUSSAINT G T. Solving geometric problems with the rotating calipers[C]//Proceedings of IEEE MELECON '83. Athens, Greece, 1983.
\bibitem{welzl} WELZL E. Smallest enclosing disks (balls and ellipsoids)[C]//MAURER H, ed. New Results and New Trends in Computer Science. Lecture Notes in Computer Science, vol. 555. Berlin, Heidelberg: Springer, 1991: 359--370. DOI: 10.1007/BFb0038202.
\bibitem{jung} JUNG H. Ueber die kleinste Kugel, die eine räumliche Figur einschliesst[J]. Journal für die reine und angewandte Mathematik, 1901, 123: 241--257. DOI: 10.1515/crll.1901.123.241.
\bibitem{tokekar} TOKEKAR P, ISLER V. Sensor placement and selection for bearing sensors with bounded uncertainty[C]//2013 IEEE International Conference on Robotics and Automation. IEEE, 2013: 2515--2520. DOI: 10.1109/ICRA.2013.6630920.
\bibitem{bertsimas2011} BERTSIMAS D, BROWN D B, CARAMANIS C. Theory and applications of robust optimization[J]. SIAM Review, 2011, 53(3): 464--501. DOI: \url{10.1137/080734510}.
\bibitem{jones1993} JONES D R, PERTTUNEN C D, STUCKMAN B E. Lipschitzian optimization without the Lipschitz constant[J]. Journal of Optimization Theory and Applications, 1993, 79(1): 157--181. DOI: \url{10.1007/BF00941892}.
\bibitem{huyer1999} HUYER W, NEUMAIER A. Global optimization by multilevel coordinate search[J]. Journal of Global Optimization, 1999, 14(4): 331--355. DOI: \url{10.1023/A:1008382309369}.
\bibitem{mavrotas2009} MAVROTAS G. Effective implementation of the $\epsilon$-constraint method in Multi-Objective Mathematical Programming problems[J]. Applied Mathematics and Computation, 2009, 213(2): 455--465. DOI: \url{10.1016/j.amc.2009.03.037}.
\bibitem{deberg2008} DE BERG M, CHEONG O, VAN KREVELD M, OVERMARS M. Computational Geometry: Algorithms and Applications[M]. 3rd ed. Berlin, Heidelberg: Springer, 2008. DOI: 10.1007/978-3-540-77974-2.
\bibitem{croes1958} CROES G A. A method for solving traveling-salesman problems[J]. Operations Research, 1958, 6(6): 791--812. DOI: 10.1287/opre.6.6.791.
\end{thebibliography}
\label{end:body}


\end{document}
